\documentclass[11pt]{article}
\usepackage[a4paper,margin=25mm]{geometry}
\usepackage{amsmath,amssymb,amsthm,microtype}
\usepackage[hidelinks]{hyperref}
\newtheorem{theorem}{Theorem}\newcommand{\E}{\mathbb E}\newcommand{\Pp}{\mathbb P}
\title{Combining all conditional counts by one crossing identity}
\author{}\date{28 October 2026}
\begin{document}\maketitle

Separate upper bounds on incoming counts discard cancellation with outgoing
counts. The following identity retains every count of the same iid law.
Its leading loss has the correct sign for all noise levels. The proof is a
finite calculation; no Poisson theorem, tail truncation, or parameter grid
enters it. The unequal binomial weights still require a separate estimate.

\begin{theorem}\label{thm:mesh}
Let $U_0,U_1$ be integrable, centered conditional-noise laws with
\[
 U_0\le a_0,\qquad -a_1\le U_1\le A_1,
 \quad a_0,a_1,A_1\ge0,\quad a_0+a_1<1,\quad A_1+a_1\le1.
\]
Take $2n$ mutually independent variables $U_{0i},U_{1i}$ with these laws,
and define
\[
 T_j=j+\sum_{i\le j}U_{1i}+\sum_{i>j}U_{0i},\quad 0\le j\le n.
\]
If $na_0\le t<n(1-a_1)$, there is a unique $K\in\{0,\ldots,n-1\}$
with $T_K\le t<T_{K+1}$. Put
\[
 b=\min\{1-A_1-a_1,1-a_0-a_1\}\ge0,\qquad R=t-T_K.
\]
Then
\begin{equation}\label{eq:gain}
              g_n(t):=t-\E K\ \ge\ \E\min\{R,b\}\ \ge0.
\end{equation}
More precisely, let $D=1+U_1$ and define the nonnegative kernel
\begin{align*}
 h(x)&=\E(x-U_0)_+-\E(x-D)_+-\Pp(D\le x),\\
 R_i&=t-(i-1)-\sum_{l<i}U_{1l}-\sum_{l>i}U_{0l}.
\end{align*}
The exact cancellation identity is
\begin{equation}\label{eq:identity}
                       g_n(t)=\sum_{i=1}^n\E h(R_i).
\end{equation}
\end{theorem}
\begin{proof}
Each increment $T_i-T_{i-1}=1+U_{1i}-U_{0i}$ is at least
$1-a_0-a_1>0$. The threshold assumptions give $T_0\le t<T_n$.
Thus the closed lower and strict upper crossing events partition the
sample space, even when $t$ is an atom of a conditional sum.

Set $m=1-a_1$. If $x<m$, both terms involving $D$ in the definition
of $h$ vanish, so
\begin{equation}\label{eq:hlow}
                         h(x)=\E(x-U_0)_+\ge0.
\end{equation}
If $x\ge m>a_0$, then $\E(x-U_0)_+=x$. Since $\E D=1$,
put--call parity and the bound $D\le1+A_1$ give
\begin{align}\label{eq:hhigh}
 h(x)&=\Pp(D>x)-\E(D-x)_+\notag\\
     &\ge(1-A_1-a_1)\Pp(D>x)\ge b\Pp(D>x).
\end{align}
This includes $x=m$ with its strict survival probability.

Let $L_j(t)=\E(t-T_j)_+$. The variables omitted from $R_i$ are
independent of it. Consequently the definition of $h$ gives
\begin{equation}\label{eq:telescoping}
 \E h(R_i)=L_{i-1}(t)-L_i(t)-\Pp(T_i\le t).
\end{equation}
Sum over $i$. The stop-loss expectations telescope, with
$L_0(t)=t-\E T_0=t$ and $L_n(t)=0$. Moreover
$\E K=\sum_{i=1}^n\Pp(T_i\le t)$. This proves
\eqref{eq:identity} and $g_n(t)\ge0$.

For the stronger age bound, condition on $R_i=x$. Crossing at $i$
means $U_{0i}\le x<D_i$, and its age is $x-U_{0i}$.
If $x<m$, then $D_i>x$ surely, and \eqref{eq:hlow} is at least the conditional
expected value of $\min\{x-U_{0i},b\}$ on this event.
If $x\ge m$, then $U_{0i}\le a_0<x$ surely, the crossing probability
is $\Pp(D_i>x)$, and its age is at least $m-a_0\ge b$.
Now \eqref{eq:hhigh} applies. Summing over the unique crossing proves
\eqref{eq:gain}. All expectations are finite, and no stopping-time or
limiting argument has been used.
\end{proof}

\section*{An exact replacement for the separate count sums}
For $Y=B+U$ with $\Pp(B=1)=p$ and these conditional laws, the conditional
distribution of a raw iid sum given $j$ successes is precisely that of
$T_j$. With $F_B$ the $\mathrm{Bin}(n,p)$ CDF,
\begin{equation}\label{eq:cdf}
            \Pp(Y_1+\cdots+Y_n\le t)=\E F_B(K).
\end{equation}
The coupling is auxiliary: it preserves each conditional law, hence
preserves the actual iid mixture. Only those conditional marginals are
needed; no claim about the original joint law of its count and $K$ is made.

For the acquired right strip, $a_0+a_1\le.072$ and $A_1=.6$, one may
take $b=.328$ in \eqref{eq:gain}; for the left strip, inward sum at most
$.068$ and $A_1=.7$ give $b=.232$. For $n\ge42$, $p\in[.25,.53]$,
$s=\E U^2$, $ns\le.616$ and $|H|\le1.84$, put
$t=np+H\sqrt{p(1-p)+s}$. Then
\[
 .25n-.948<t<.53n+.948,\qquad
                   na_0\le t<n(1-a_1).
\]
Indeed $\sqrt{1/4+.616/42}<.515$ and $1.84(.515)<.948$;
the remaining two inequalities hold already at $n=42$. Thus the
cancellation identity covers the whole retained acquired noise range.
It makes no assertion that the acquisition hypotheses themselves hold
without their existing contact and predecessor dependencies.

\section*{The weighted identity isolates the remaining difficulty}
Write $\pi_j=\Pp(B_n=j)$, $L_j(t)=\E(t-T_j)_+$,
$B_j(t)=L_j(t)-(t-j)_+\ge0$ (by Jensen), and
$g_i(t)=\E h(R_i)\ge0$. Multiply \eqref{eq:telescoping} by $\pi_i$ and sum by parts.
Because the count-zero CDF is one and $L_0=t$, $L_n=0$, this gives
\begin{equation}\label{eq:parts}
 \Pp(Y_1+\cdots+Y_n\le t)
 =\pi_0+\pi_1t+\sum_{j=1}^{n-1}(\pi_{j+1}-\pi_j)L_j(t)
                 -\sum_{i=1}^n\pi_i g_i(t).
\end{equation}
Let $k=\lfloor t\rfloor$ and $u=t-k\in[0,1)$. Replacing $L_j$ by
$(t-j)_+$ in the first three terms gives the linear interpolation
$F_B(k)+u\pi_{k+1}$. Hence the exact formula is
\begin{equation}\label{eq:weighted}
 \Pp(Y_1+\cdots+Y_n\le t)-F_B(k)
 =u\pi_{k+1}+\sum_{j=1}^{n-1}(\pi_{j+1}-\pi_j)B_j(t)
                 -\sum_{i=1}^n\pi_i g_i(t).
\end{equation}
Both $B_j$ and $g_i$ are nonnegative. Only the increasing side of the
binomial mass function contributes unfavorably to the curvature sum.
No separate incoming-count tail appears.

For the following normalized comparison assume additionally $0<p<1$
and $\E|U_0|^3+\E|U_1|^3<\infty$, as holds in the acquired bounded
strips. Put $v=p(1-p)$, $s=\E U^2$,
$\rho=\E|Y-p|^3/(v+s)^{3/2}$ and $\rho_B=(1-2v)/\sqrt v$.
For the direct comparison at this integer anchor,
\begin{align}\label{eq:delta}
 \Delta_C={}&-\sqrt n\sum_{i=1}^n\pi_i g_i(t)
       +\sqrt n\sum_{j=1}^{n-1}(\pi_{j+1}-\pi_j)B_j(t)
       +\sqrt n\,u\pi_{k+1}\notag\\
 &-\sqrt n\left[
 \Phi\!\left(\frac{t-np}{\sqrt{n(v+s)}}\right)
 -\Phi\!\left(\frac{k-np}{\sqrt{nv}}\right)\right]
                  +C(\rho_B-\rho).
\end{align}
Every identity is exact at the actual sample size. The remaining task is
to absorb the binomial curvature and the phase, normal and cubic repairs
in the common nonnegative loss. A sign for the unweighted loss alone does
not do this. Therefore the current general count certificates remain
ancestors of the accepted comparison until that joint inequality is proved.

For comparison with a direct stopped-sum calculation, put $W_j=T_j-j$.
Its selected mean is generally nonzero:
\[
 \E W_{K+1}=-\sum_{i=1}^n
   \E\!\left[U_{0i}\mathbf1_{\{T_{i-1}\le t\}}\right]\ge0.
\]
The sign follows by conditioning on the omitted-pair sum and centering
$U_{0i}$. The telescoping proof includes this bias automatically; setting
it to zero would be an unjustified use of optional stopping.

\section*{Controls and limitations}
The companion exact checks reconstruct all conditional sums and the local
kernel $h$ separately, including closed collisions and strict crossings.
They check \eqref{eq:identity}, \eqref{eq:gain}, \eqref{eq:cdf}, \eqref{eq:weighted} and the
nonzero selected-noise bias. They also include a false extension: take
$U_0=0$ and $D=9/10,39/20$ with probabilities $19/21,2/21$.
At $n=2$, $t=9/10$, $\E K-t=1/210>0$; here $A_1+a_1=21/20>1$.
No assertion of strict gain proportional to $s$, or of the full sharp
Berry--Esseen bound, follows from this note alone.
\end{document}
